\kchapter{Ingress, Gateway API, and Network Policies}{ing}

\begin{objectives}
\begin{itemize}
  \item how HTTP traffic enters a cluster through one entry point;
  \item the Ingress resource and the newer, more powerful Gateway API;
  \item TLS certificates with cert-manager;
  \item NetworkPolicies: a firewall between Pods;
  \item what a CNI plugin is and how to choose one.
\end{itemize}
\end{objectives}

\section{The problem: many services, one front door}

You have a frontend, an API, and an admin panel. You want:

\begin{itemize}
  \item \kc{shop.example.com} $\rightarrow$ frontend;
  \item \kc{shop.example.com/api} $\rightarrow$ API;
  \item \kc{admin.example.com} $\rightarrow$ admin;
  \item HTTPS everywhere, with one public IP and one load balancer.
\end{itemize}

This is \kterm{Layer-7} (HTTP) routing. Kubernetes models it with two APIs:
the original \textbf{Ingress} and its successor, the \textbf{Gateway API}.
Both are only \emph{descriptions}; a \kterm{controller} (a proxy such as
Envoy, NGINX, HAProxy, or Traefik running in your cluster) reads them and
does the actual routing.

\begin{center}
\resizebox{\linewidth}{!}{%
\begin{tikzpicture}[
  b/.style={rounded corners=5pt, minimum height=1.1cm, align=center, font=\scriptsize\bfseries},
  >={Stealth[length=2.6mm]}]
  \node[b, fill=k8sgray!20, minimum width=2cm] (u) at (0,0) {\faGlobe\\Internet};
  \node[b, fill=k8sorange, text=white, minimum width=2.4cm] (lb) at (3.2,0) {Cloud LB\\[-1pt]\mdseries one public IP};
  \node[b, fill=k8sblue, text=white, minimum width=3.0cm, minimum height=1.6cm] (g) at (7.0,0)
       {\kimg[0.6cm]{ing}\\Gateway / Ingress\\controller (Envoy\ldots)};
  \node[b, fill=k8slight, draw=k8sblue, minimum width=3.1cm] (s1) at (11.6,1.6) {\kicon{svc}\ frontend\\[-1pt]\mdseries\ttfamily shop.example.com/};
  \node[b, fill=k8slight, draw=k8sblue, minimum width=3.1cm] (s2) at (11.6,0)   {\kicon{svc}\ api\\[-1pt]\mdseries\ttfamily shop.example.com/api};
  \node[b, fill=k8slight, draw=k8sblue, minimum width=3.1cm] (s3) at (11.6,-1.6){\kicon{svc}\ admin\\[-1pt]\mdseries\ttfamily admin.example.com};
  \draw[->, thick] (u) -- node[above, font=\tiny]{HTTPS} (lb);
  \draw[->, thick] (lb) -- (g);
  \foreach \s in {s1,s2,s3} \draw[->, thick, k8sblue] (g) -- (\s.west);
\end{tikzpicture}}
\end{center}

\section{Ingress}

\begin{yamlcode}{shop-ingress.yaml}
apiVersion: networking.k8s.io/v1
kind: Ingress
metadata:
  name: shop
spec:
  ingressClassName: traefik        # which controller should handle this
  tls:
    - hosts: [shop.example.com]
      secretName: shop-tls          # certificate + key stored in a Secret
  rules:
    - host: shop.example.com
      http:
        paths:
          - path: /api
            pathType: Prefix
            backend:
              service: { name: api, port: { number: 80 } }
          - path: /
            pathType: Prefix
            backend:
              service: { name: frontend, port: { number: 80 } }
\end{yamlcode}

Ingress is stable and widely supported, but it is limited: only host/path
routing is standard, and everything else (rewrites, timeouts, canaries, header
matching) is done with controller-specific annotations that are not portable.

\begin{warningbox}
The popular community \textbf{ingress-nginx} controller has been
\textbf{retired} by the Kubernetes project (best-effort maintenance ended in
March 2026). Existing clusters should plan a migration. For new work, prefer the
\textbf{Gateway API} with a maintained implementation such as Envoy Gateway,
Istio, Cilium, NGINX Gateway Fabric, Traefik, or your cloud provider's gateway.
\end{warningbox}

\section{Gateway API}

The \kterm{Gateway API} is the modern, role-oriented, expressive replacement
for Ingress (GA since v1.0 in 2023). It splits responsibility across three
kinds of objects:

\begin{center}
\ktablestart
\begin{tabularx}{\linewidth}{P{2.7cm} P{3.4cm} L}
\kheader \kth{Object} & \kth{Owned by} & \kth{Describes} \\
\kc{GatewayClass} & Infrastructure provider & Which controller implements gateways (like a StorageClass for traffic). \\
\kc{Gateway} & Platform / cluster operator & A concrete entry point: listeners, ports, hostnames, TLS. \\
\kc{HTTPRoute} (+ \kc{GRPCRoute}, \kc{TLSRoute}\ldots) & Application team & How requests are matched and forwarded to Services. \\
\end{tabularx}
\end{center}

\begin{yamlcode}{gateway.yaml}
apiVersion: gateway.networking.k8s.io/v1
kind: Gateway
metadata:
  name: public
  namespace: infra
spec:
  gatewayClassName: eg                 # e.g. Envoy Gateway
  listeners:
    - name: https
      protocol: HTTPS
      port: 443
      hostname: "*.example.com"
      tls:
        mode: Terminate
        certificateRefs:
          - name: wildcard-example-com
      allowedRoutes:
        namespaces:
          from: Selector
          selector:
            matchLabels:
              gateway-access: "true"
\end{yamlcode}

\begin{yamlcode}{shop-route.yaml}
apiVersion: gateway.networking.k8s.io/v1
kind: HTTPRoute
metadata:
  name: shop
  namespace: shop
spec:
  parentRefs:
    - name: public
      namespace: infra
  hostnames: ["shop.example.com"]
  rules:
    - matches:
        - path: { type: PathPrefix, value: /api }
      backendRefs:
        - name: api-v1
          port: 80
          weight: 90                   # canary: 90% to v1 ...
        - name: api-v2
          port: 80
          weight: 10                   # ... and 10% to v2
    - matches:
        - headers:
            - name: x-beta-user
              value: "true"
      backendRefs:
        - name: frontend-beta
          port: 80
    - backendRefs:                     # default route
        - name: frontend
          port: 80
\end{yamlcode}

Weighted backends, header matching, redirects, rewrites, timeouts, and request
mirroring are all part of the \emph{standard} API --- no annotations needed.

\section{TLS certificates with cert-manager}

\kicon{cert-manager}~\textbf{cert-manager} automates certificates from Let's
Encrypt (or your internal CA) and stores them as Secrets that Ingress and
Gateways reference. It renews them before expiry.

\begin{yamlcode}{cluster-issuer.yaml}
apiVersion: cert-manager.io/v1
kind: ClusterIssuer
metadata:
  name: letsencrypt
spec:
  acme:
    server: https://acme-v02.api.letsencrypt.org/directory
    email: ops@example.com
    privateKeySecretRef:
      name: letsencrypt-account-key
    solvers:
      - http01:
          gatewayHTTPRoute:
            parentRefs:
              - name: public
                namespace: infra
\end{yamlcode}

\section{Network Policies: a firewall for Pods}

By default \textbf{every Pod can talk to every other Pod}. A
\kterm{NetworkPolicy} restricts that. Once any policy selects a Pod for a
direction (ingress or egress), only traffic explicitly allowed by some policy
is permitted in that direction.

\begin{keyidea}
NetworkPolicies are \textbf{allow-lists} and they are \textbf{additive}. Start
with ``deny all'' in a namespace, then add policies that open exactly the paths
your application needs.
\end{keyidea}

\begin{yamlcode}{default-deny.yaml}
apiVersion: networking.k8s.io/v1
kind: NetworkPolicy
metadata:
  name: default-deny
  namespace: shop
spec:
  podSelector: {}           # every Pod in the namespace
  policyTypes: [Ingress, Egress]
\end{yamlcode}

\begin{yamlcode}{allow-api-to-db.yaml}
apiVersion: networking.k8s.io/v1
kind: NetworkPolicy
metadata:
  name: db-allow-api
  namespace: shop
spec:
  podSelector:
    matchLabels: { app: db }           # this policy protects the DB Pods
  policyTypes: [Ingress]
  ingress:
    - from:
        - podSelector:
            matchLabels: { app: api }  # only API Pods ...
      ports:
        - port: 5432                    # ... on the Postgres port
          protocol: TCP
---
apiVersion: networking.k8s.io/v1
kind: NetworkPolicy
metadata:
  name: allow-dns-egress
  namespace: shop
spec:
  podSelector: {}
  policyTypes: [Egress]
  egress:
    - to:
        - namespaceSelector:
            matchLabels:
              kubernetes.io/metadata.name: kube-system
          podSelector:
            matchLabels:
              k8s-app: kube-dns
      ports:
        - { port: 53, protocol: UDP }
        - { port: 53, protocol: TCP }
\end{yamlcode}

\begin{warningbox}
After a default-deny on \emph{egress}, Pods cannot even resolve DNS. Always add
the DNS egress rule above. Also note the YAML subtlety: \kc{namespaceSelector}
and \kc{podSelector} in the \emph{same} list item mean AND; as \emph{separate}
items (each with its own dash) they mean OR.
\end{warningbox}

\begin{warningbox}
\textbf{NetworkPolicies need a CNI plugin that enforces them.} Calico and Cilium
do; Flannel alone does not. If your plugin ignores them, your policies are
silently decorative.
\end{warningbox}

\section{CNI: the network plugin}

The \kterm{Container Network Interface} (CNI) is the plugin standard that gives
Pods their IPs and connects them. Popular choices:

\begin{center}
\ktablestart
\begin{tabularx}{\linewidth}{P{2.4cm} L}
\kheader \kth{Plugin} & \kth{Highlights} \\
Cilium & eBPF-based, very fast, NetworkPolicies at L3--L7, can replace
         kube-proxy, rich observability (Hubble). \\
Calico & Mature, BGP routing, strong NetworkPolicy support, widely used on-prem. \\
Flannel & Simple overlay network; no policy enforcement by itself. \\
Cloud CNIs & AWS VPC CNI, Azure CNI, GKE Dataplane V2 --- Pods get VPC-native IPs. \\
\end{tabularx}
\end{center}

\begin{recap}
\begin{itemize}
  \item One controller + one load balancer can route all HTTP traffic by host and path.
  \item Ingress is simple and stable; the Gateway API is the expressive, portable successor.
  \item cert-manager automates TLS certificates.
  \item NetworkPolicies are additive allow-lists; start from default-deny and allow DNS.
  \item Policies only work if your CNI enforces them.
\end{itemize}
\end{recap}
